\documentclass[11pt,a4paper]{article}
\usepackage[utf8]{inputenc}
\usepackage{mathtools, amssymb}
\usepackage{graphicx}
\graphicspath{{figures/}{data/drafts/figures/}{../}{data/}}
\usepackage{booktabs}
\usepackage{natbib}
\usepackage[margin=1in]{geometry}
\usepackage{physics}
\usepackage{xcolor}
\usepackage{authblk}
\usepackage{siunitx}
\usepackage{hyperref}

\title{Metric-Consistent Spectral Finite Element Decomposition of Stable Boundary Layer Regimes: \\
A Riemannian Manifold Approach to Sub-Mesoscale Wave--Turbulence Partitioning}

\author{David E. England\textsuperscript{1,*}\thanks{ORCID: \href{<https://orcid.org/0009-0001-2095-6646}{0009-0001-2095-6646}}}>
\affil{\textsuperscript{1}The University of Alabama in Huntsville, 301 Sparkman Drive, Huntsville, AL 35899, USA}
\affil{\textsuperscript{*}Corresponding author: <David.England@UAH.Edu>}

\date{\today}

\begin{document}

\maketitle

\begin{abstract}
Nocturnal stable boundary layers (SBLs) exhibit complex interactions between three dominant scales: large-scale mesoscale pressure forcing, horizontally propagating internal gravity waves (IGWs), and highly localized turbulent shear bursts. Traditional diagnostic frameworks conflate these regimes, producing high variance in inversion-layer flux estimates and inconsistent model-observation comparisons. We present a \emph{metric-consistent Riemannian pseudospectral} framework that projects sparse tower observations onto a non-uniform computational grid designed explicitly to maximize sensitivity near the surface (\SI{2.85}{\milli\meter} effective node resolution at $z = \SI{1.5}{\meter}$) while maintaining numerical stability at upper boundaries ($\kappa(\mathbf{M}) = 4619$). Using Chebyshev polynomials as the spectral basis and solving the rank-deficient least-squares problem via thin singular value decomposition (SVD), we extract three physically decoupled diagnostic features: effective modal dimension $D_{\mathrm{eff}}$, wave energy fraction $F_W$, and spectral curvature $\chi_N$.

Applying this framework to the CASES-99 campaign stable window (22--31 October 1999), we identify \emph{three statistically separable regimes} with a silhouette clustering score of $\bar{S} = 0.582$ for $K=3$: (1) Continuous Turbulence ($D_{\mathrm{eff}} > 20$, $F_W < 0.15$), (2) Wave-Dominated ($D_{\mathrm{eff}} \sim 4$--$6$, $F_W > 0.75$), and (3) Intermittent Shear Bursts ($2 < D_{\mathrm{eff}} < 15$, $0.15 < F_W < 0.50$). Crucially, the regime-specific rank compression ($\Delta_{\mathrm{rank}} = 24$) proves that atmospheric dynamics naturally occupy lower-dimensional manifolds and are not mathematical artifacts of basis truncation.

We extend the theoretical framework by establishing an explicit physical link to the Ozmidov buoyancy scale $L_O$, resolving the grid-topology mismatch via a generalized mathematical proof-of-concept for a ``Virtual Tower'' interpolation operator ($\mathcal{P}_{\text{VT}}$) mapping numerical model grids ($\eta$-coordinates) directly onto the computational manifold, and formulating a $2/3$-rule spectral anti-aliasing filter for forward-predictive implementations. A synthetic gravity-wave reflection test confirms that the spectral partition functions suppress upper-boundary reflection by approximately $3.3\times 10^4$, validating the sponge-layer boundary conditions.
\end{abstract}

\noindent \textbf{Keywords:} stable boundary layers, spectral methods, manifold learning, gravity waves, Ozmidov scale, virtual tower interpolation

\newpage
\tableofcontents
\newpage

% ============================================================================
\section{Introduction}
\label{sec:intro}

Nocturnal stable boundary layers (SBLs) remain among the most poorly understood and inconsistently modeled features of the lower atmosphere \citep{Mahrt2014, Acevedo2016}. Unlike daytime convective boundary layers, where surface heating drives turbulence vertically, nocturnal SBLs exhibit a fundamental instability: strong temperature inversions suppress vertical mixing, yet persistent wind shear continues to generate turbulent kinetic energy (TKE) locally. The result is an environment dominated by \emph{intermittency}---alternating periods of quiescent, laminar flow interrupted by sudden shear-driven turbulent bursts that collapse and dissipate over minutes to tens of minutes \citep{Mahrt1999, Terradellas2001}.

Superimposed on this intermittency are horizontally propagating internal gravity waves (IGWs), generated by drainage flows over complex terrain or by baroclinic vorticity forcing at mesoscale boundaries \citep{Vosper2006, Sun2015}. As synthesized by \citet{Sun2015}, stably stratified environments naturally support an overlapping continuum of anisotropic, intermittent turbulence and wavelike motions. These waves transport energy horizontally and can locally steepen density gradients, triggering instability and turbulent collapse. However, traditional single-column decomposition schemes (e.g., Richardson number thresholds, wavelet filtering) treat waves and turbulence as completely separable, missing the \emph{coupled} dynamics that govern real boundary-layer transitions.

\begin{figure}[htbp]
    \centering
    \includegraphics[width=0.75\linewidth, angle=-90]{figures/Night_Sky_Turbulence_Classification_Framework.png.pdf}
    \caption{Classification framework for night sky stable boundary layer regimes mapping scale-dependent multi-regime transitions.}
    \label{fig:framework}
\end{figure}

\subsection{The MOST Crisis}

The foundational framework for stable boundary-layer parameterization remains Monin--Obukhov similarity theory (MOST) \citep{Obukhov1946, Paulson1970}, which predicts that normalized gradients and fluxes depend only on a single stability parameter:
\begin{equation}
\zeta = \frac{z}{L} = \frac{z \kappa g}{\theta_{\mathrm{ref}}} \frac{\overline{T'w'}}{u_*^2 \theta_{\mathrm{ref}}}
\label{eq:zeta}
\end{equation}
where $L$ is the Obukhov length scale, $u_*$ is friction velocity, and $\overline{T'w'}$ is the sensible heat flux. For stable conditions ($\zeta > 0$), MOST predicts:
\begin{equation}
\phi_m(\zeta) = 1 + \beta_m\zeta, \quad \phi_h(\zeta) = 1 + \beta_h\zeta,
\end{equation}
with universal constants $\beta_m \approx 5$ and $\beta_h \approx 5$ (though empirical estimates vary widely, $\beta_m \in [3, 10]$, $\beta_h \in [3, 15]$ \citep{Dyer1974, Högström1988}).

However, observational datasets from tower campaigns---CASES-99 \citep{Poulos2002}, SHEBA \citep{Uttal2002}, FINN \citep{Mahrt2013}---reveal systematic deviations from MOST scaling. Flux-profile relationships show distinct hysteresis loops: the same gradient Richardson number yields entirely different flux estimates depending on whether the system is transitioning into or out of an inversion layer. Furthermore, wave-affected periods show scaling parameters values far outside the universal range, with non-monotonic behavior in $\zeta$, while intermittent turbulence produces rapid switches in effective $L$ that render time-averaged diagnostics fundamentally ambiguous.

As established in the comprehensive review by \citet{Sun2015}, state-of-the-art numerical weather prediction models fail systematically to simulate this seemingly random turbulence intensification because they lack a rigorous mechanism to parameterize wave-turbulence interactions within stably stratified geophysical flows.

\subsection{Spectral Decomposition Approaches}

Recent efforts have attempted to isolate waves from turbulence using wavelet filtering \citep{Vosper2006, Sun2015} or empirical mode decomposition (EMD) \citep{Huang1998}. While wavelet filtering decomposes signals into time-frequency atoms, the basis choice (Morlet, Daubechies) remains fundamentally arbitrary and introduces non-local Gibbs ringing near physical discontinuities. EMD extracts intrinsic mode functions adaptively, but its outputs are non-orthogonal, and the decomposition lacks mathematical uniqueness. Similarly, Proper Orthogonal Decomposition (POD) and Dynamic Mode Decomposition (DMD) techniques \citep{Schmid2010} extract dominant modes cleanly but require dense spatiotemporal snapshot matrices, which are entirely unavailable from sparse vertical instrument masts.

None of these methods preserve the physical constraint that projections must respect the \emph{natural metric} of the atmospheric system. A wave acting on a stable density gradient operates under a different effective mass than a turbulent eddy in the same layer. Standard spectral projections treat both fields identically, obscuring the primary coupling modes highlighted by \citet{Sun2015}.

\subsection{Contribution: Metric-Consistent Manifold Decomposition}

This paper presents a fundamentally different approach: instead of projecting data onto arbitrary basis functions, we construct a \emph{Riemannian metric} $\mathbf{M}$ that weights projections according to physical density and stratification. The metric ensures that wave energy is cleanly localized to low-frequency Chebyshev modes that respect the stable density gradient, while turbulent kinetic energy is concentrated in high-frequency modes highly sensitive to localized vertical wind shear.

Key innovations presented herein include:
\begin{enumerate}
    \item \textbf{Hyperbolic compactification} of the vertical coordinate: Non-uniform grid packing ($\Delta z_{\min} = \SI{2.85}{\milli\meter}$ at the surface, $\Delta z_{\max} = \SI{9.43}{\meter}$ aloft) captures acute $\SI{5}{\kelvin\per\meter}$ inversions without wasting degrees of freedom on upper air layers.
    \item \textbf{Thin SVD-truncated matrix projection}: With 8 instrument levels but an $N=32$ order spectral expansion, rank deficiency is structurally unavoidable. We solve this inverse problem via an economy-size, thin Singular Value Decomposition (SVD), restricting matrix operations entirely to the data-constrained subspace.
    \item \textbf{Partition-of-unity spectral masking} $(\psi_M, \psi_W, \psi_T)$: This framework explicitly separates mesoscale mean flow, gravity-wave scales, and turbulent dissipation without introducing ad-hoc cutoff boundaries.
\end{enumerate}

We formalize the physical link between our information-theoretic metric ($D_{\mathrm{eff}}$) and the Ozmidov buoyancy scale, provide a generalized ``Virtual Tower'' interpolation framework to resolve the grid-topology mismatch for global model validation, and supply a spectral anti-aliasing scheme necessary for forward-predictive closure models.

\subsection{Structural and Theoretical Map of Sources}
\label{sec:intro:source_map}

The theoretical arc of this study is organized around three linked source families.

\paragraph{(1) Stable boundary-layer dynamics and the wave--turbulence crisis.}
The physical motivation follows directly from synthesis and field evidence showing that nocturnal stable boundary layers are governed by overlapping, anisotropic wave--turbulence bands rather than separable single-column states \citep{Sun2015, Acevedo2016}. Foundational intermittency observations \citep{Mahrt1999, Terradellas2001} establish why abrupt shear-burst transitions challenge conventional Fourier/wavelet assumptions, while canonical MOST flux-profile closuresI’ll integrate this source map into your current manuscript in v9.tex as a compact “Structural and Theoretical Map” subsection, aligned with your existing citation style and equations.

Searched for regex `Sun2015|Acevedo2016|Mahrt1999|Terradellas2001|Paulson1970|Dyer1974|Högström1988|Trefethen2000|Boyd2001|Fornberg1996|Huang1998|Schmid2010|Vosper2006|Vosper2013|Sheridan2010|Wood2014|Halliday2018` (`**/drafts/v9.tex`), 16 results

Read [](file:///Users/davidengland/Documents/GitHub/SpectralBL/drafts/v9.tex#1-1), lines 1 to 260

I’m going to add a compact “Structural and Theoretical Map of Sources” subsection in the Introduction of v9.tex, organized by the three source pillars you outlined and tied directly to your existing framing.

Generating patch (21 lines) in [](file:///Users/davidengland/Documents/GitHub/SpectralBL/drafts/v9.tex)


```

```


Checked [](file:///Users/davidengland/Documents/GitHub/SpectralBL/drafts/v9.tex), no problems found

Added your requested source-structure section into the manuscript and aligned it with your existing narrative flow.

What I changed

- Inserted a new subsection:
  - v9.tex
  - Title: Structural and Theoretical Map of Sources
- Organized it into the same three pillars you specified:
  1. Stable boundary-layer dynamics and wave–turbulence crisis
  2. Pseudospectral and numerical foundations
  3. Mesoscale model interfacing via the mathcal P_{VT} pathway
- Preserved and integrated the citation set you listed:
  - Sun2015, Acevedo2016, Mahrt1999, Terradellas2001, Paulson1970, Dyer1974, Högström1988
  - Trefethen2000, Boyd2001, Fornberg1996, Huang1998, Schmid2010
  - Vosper2006, Vosper2013, Sheridan2010, Wood2014, Halliday2018

Validation

- Checked for syntax/problems after edit: no LaTeX file-level errors reported for v9.tex.

If you want, next I can also tighten the wording in the Methods section so the new subsection and the existing equations/notation (especially the SVD paragraph) are fully style-consistent line-by-line.

Made changes. \citep{Paulson1970, Dyer1974, H\"ogstr\"om1988} provide the empirical baseline that this manifold framework is designed to supersede in strongly stratified, non-stationary windows.

\paragraph{(2) Pseudospectral and numerical foundations.}
The numerical architecture is grounded in non-periodic spectral methods, with Chebyshev $T_n$ bases and differentiation/mass constructions following standard pseudospectral practice \citep{Trefethen2000, Boyd2001, Fornberg1996}. In parallel, prior decomposition alternatives motivate the low-rank inverse design: EMD lacks orthogonal uniqueness \citep{Huang1998}, and DMD/POD-style modal extraction is snapshot-hungry under sparse tower sampling \citep{Schmid2010}. This jointly motivates the SVD-truncated operator $\mathbf{H}$ used here for physically constrained reconstruction from limited vertical levels.

\paragraph{(3) Mesoscale model interfacing and the $\mathcal{P}_{\mathrm{VT}}$ pathway.}
The ``Virtual Tower'' mapping direction is supported by studies showing strong near-surface sensitivity of stable-flow simulation skill to vertical grid structure and terrain-coupled cold-pool dynamics \citep{Vosper2006, Vosper2013, Sheridan2010}. This supports explicit inverse stretching $\xi(z)$ when projecting coarse model elevations $\mathbf{z}_{\eta}$ onto the manifold state space. Forward-operational relevance is reinforced by non-hydrostatic, mass-conserving and gravity-wave-resolving mesoscale formulations \citep{Wood2014, Halliday2018}, which align with our use of geometry-aware diagnostics for regime-aware closure and validation.

% ============================================================================
\section{Data: CASES-99 Stable Window}
\label{sec:data}

The Cooperative Atmosphere-Surface Exchange Study--1999 (CASES-99) field observational program took place from 1 October to 31 October 1999 centered near Leon, Kansas ($37.69^\circ\text{N}$, $97.64^\circ\text{W}$) \citep{Poulos2002}. The terrain at the Leon site consists of gently rolling grassland with sparse trees, a displacement height $d \approx \SI{0.1}{\meter}$, and an aerodynamic roughness length $z_0 \approx \SI{0.01}{\meter}$.

\subsection{Main Tower Array and Hardware-Specific Quality Assurance}
We use data from the main豪 CASES-99 tower array, which hosted a central \SI{55}{\meter} to \SI{60}{\meter} main mast instrumented across eight distinct vertical levels: 1.5, 5.0, 10.0, 20.0, 30.0, 40.0, 50.0, and \SI{55.0}{\meter}. To reconstruct a continuous vertical timeline without introducing instrument-induced artifacting into the pseudospectral projection, the data ingestion engine applies a multi-brand hardware quality gate to the 20~Hz resampled sonic anemometer streams. Let $\mathbf{u}_z(t)$ denote the raw velocity vector at a given instrument height $z$. The quality-control operator $\mathcal{Q}$ isolates instrument status metrics across the distinct anemometer types distributed along the mast:
\begin{equation}
\mathcal{Q}\left(\mathbf{u}_z(t)\right) =
\begin{cases}
\mathbf{NaN}, & \text{if } z \in \{1.5, 5, 30, 50\}\,\text{m and } \mathtt{diag}_z \neq 0 \\
\mathbf{NaN}, & \text{if } z \in \{10, 20, 40, 55\}\,\text{m and } \mathtt{usamples}_z < 10 \\
\mathbf{u}_z(t), & \text{otherwise}
\end{cases}
\end{equation}
where $\mathtt{diag}_z$ represents the bit-field diagnostic flag native to the Campbell Scientific CSAT3 sonic anemometers, used primarily to intercept transducer blockage from dew or nocturnal frost formation. Conversely, $\mathtt{usamples}_z$ monitors the internal sub-sampling density (out of a nominal 200~Hz base rate averaged down to 20~Hz) for the Applied Technologies, Inc. (ATI) sonic anemometers. Sub-sampling drops below the threshold ($\mathtt{usamples}_z < 10$) indicate intermittent data frame dropouts induced by high-frequency sub-mesoscale shear bursts. Independent temperature cross-checks are maintained via collocated Omega Type-E thermocouples across all eight evaluation tiers.

\subsection{Analysis Window: October 22--31 Stable Period}

We focus on the intensive observation period (IOP) from 22--31 October 1999, which was characterized by strong radiative cooling and high-amplitude stable inversions. This 10-day window provides consistent dry conditions (dew-point depression $> \SI{15}{\kelvin}$) minimizing latent heat effects, weak large-scale subsidence ($\sim \SI{2}{\kelvin\per\text{day}}$ synoptic cooling) allowing fine-scale stable dynamics to dominate, and strong, coherent gravity-wave activity documented across multiple independent diagnostic tracking systems \citep{Poulos2002}. The analysis encompasses \emph{1440 10-minute-averaged profiles} ($n_t = 144$ hours) with zero missing values at the eight tower levels after QA filtering.

\subsection{Data Quality Assurance and Verification Metrics}

Raw time series profiles were subjected to threshold rejection ($\theta < \SI{260}{\kelvin}$, $u > \SI{25}{\meter\per\second}$), vertical monotonicity verification ($\partial \theta / \partial z > -\SI{0.01}{\kelvin\per\meter}$), and rank-sufficiency checking. This quality screening pipeline retained 1386 valid profiles ($96.3\%$ of the raw data).

To mathematically justify the use of a Gaussian Mixture Model (GMM) for multi-regime classification without necessitating a prior, information-lossy Principal Component Analysis (PCA) projection, we mapped the global linear independence of the feature array. Table \ref{tab:correlation_matrix} details the computed cross-feature Pearson correlation profile ($R$) across the entire dataset footprint.

\begin{table}[h]
\centering
\caption{Global Feature Pearson Correlation Matrix ($R$) Across the CASES-99 Campaign Footprint}
\label{tab:correlation_matrix}
\begin{tabular}{l|cccc}
\hline
Metric Feature Array & $D_{\mathrm{eff}}$ & $F_W$ & $\chi_N$ & $\mathrm{Ri}_g$ \\
\hline
$D_{\mathrm{eff}}$   & 1.00 & -0.01 &  0.02 &  0.01 \\
$F_W$                & -0.01 &  1.00 & -0.49 & -0.00 \\
$\chi_N$             &  0.02 & -0.49 &  1.00 &  0.00 \\
$\mathrm{Ri}_g$      &  0.01 & -0.00 &  0.00 &  1.00 \\
\hline
\end{tabular}
\end{table}

The empirical correlation bounds confirm that no feature pair approaches the critical collinearity threshold ($|R| > 0.80$). The near-zero correlation between local gradient stability ($\mathrm{Ri}_g$) and the structural curvature shape parameter ($\chi_N$) proves that macro-geometric profiling captures separate boundary layer physics completely invisible to localized gradient formulations.

% ============================================================================
\section{Methods}
\label{sec:methods}

\subsection{Unified Manifold Workspace and Coordinate Mapping}
\label{sec:methods:manifold}

Standard pseudospectral frameworks relying on uniform or Chebyshev-cosine grids distribute computational degrees of freedom symmetrically across the spatial domain. Such configurations are fundamentally unoptimized for the acute vertical gradients characterizing the nocturnal stable boundary layer (SBL), where sub-mesoscale gravity waves and shear instabilities concentrate within the lowest meters of the surface layer ($1.5 \le z \le 5.0\,\text{m}$). To preferentially focus node density where structural geometric derivatives are maximized, we establish a \emph{UnifiedManifoldWorkspace} utilizing a fractional hyperbolic coordinate stretching function $\mathcal{T}(\xi) \rightarrow z$. This operator transforms the bounded computational space $\xi \in [-1, 1]$ into the physical elevation domain defined by the tower instruments ($z_{\min} = 1.5\,\text{m}$ to $z_{\max} = 55.0\,\text{m}$):
\begin{equation}
z(\xi) = z_{\min} + \frac{z_{\max} - z_{\min}}{1 - \alpha_s^2} \left[ \frac{(1+\alpha_s)(1+\xi)}{1 + \alpha_s(2+\xi)} \right]
\label{eq:hyperbolic_stretch}
\end{equation}
where $\alpha_s \in (0, 1)$ represents the non-dimensional grid compression parameter. The analytical transformation Jacobian resulting from this compactification is defined as:
\begin{equation}
J(\xi) = \frac{dz}{d\xi} = \frac{(z_{\max} - z_{\min})(1 + \alpha_s)^2}{(1 - \alpha_s^2) \left[1 + \alpha_s(2+\xi)\right]^2}
\label{eq:jacobian_stretched}
\end{equation}

Fixing the operational baseline at $\alpha_s = 0.05$ compresses the local node spacing to an ultra-dense scale ($\Delta z_{\min} = J(-1) \cdot \Delta \xi_{\min} \approx 2.85\,\text{mm}$) directly above the surface boundary, while allowing node separation to expand monotonically toward the top boundary ($\Delta z_{\max} \approx 9.43\,\text{m}$ as $z \rightarrow 55\,\text{m}$). To justify this spatial geometry against alternative topologies, a rigorous sensitivity study was performed across a parameterized continuum of stretching bounds (Table~\ref{tab:stretch_sensitivity}).

\begin{table}[h]
\centering
\caption{Mass Matrix Conditioning and Grid Resolution Sensitivity Across Parameterized Stretching Bounds}
\label{tab:stretch_sensitivity}
\begin{tabular}{c|rrrc}
\hline
$\alpha_s$ & $\Delta z_{\min}$ (at $1.5\,\text{m}$) & $\Delta z_{\max}$ (at $55\,\text{m}$) & $\kappa(\mathbf{M})$ & Operational Status \\
\hline
0.01 & $0.11\,\text{mm}$ & $43.20\,\text{m}$ & $1.21 \times 10^5$ & Ill-Conditioned Boundary \\
0.05 & $\mathbf{2.85\,\text{mm}}$ & $\mathbf{9.43\,\text{m}}$ & $\mathbf{4619}$ & \textbf{Optimized Baseline} \\
0.10 & $10.62\,\text{mm}$ & $4.85\,\text{m}$ & $924$ & Suppressed Surface Sensitivity \\
0.50 & $98.40\,\text{mm}$ & $1.82\,\text{m}$ & $63$ & Quasi-Uniform Limit \\
\hline
\end{tabular}
\end{table}

The computational manifold embeds this non-uniform spacing into the discrete physics via a metric-weighted mass matrix $\mathbf{M} \in \mathbb{R}^{(N+1) \times (N+1)}$. This structure discretizes the weighted continuous inner product space using high-order Gauss--Lobatto quadrature across $K_q = 72$ nodes:
\begin{equation}
M_{mn} = \int_{-1}^{1} T_m(\xi) T_n(\xi) J(\xi) \, d\xi \approx \sum_{k=1}^{K_q} w_k \, T_m(\xi_k) T_n(\xi_k) J(\xi_k)
\label{eq:mass_matrix_quad}
\end{equation}
where $T_n(\xi)$ denotes the $n$-th Chebyshev polynomial of the first kind. At $\alpha_s = 0.05$, the metric-weighted system maintains a stable spectral profile ($\kappa(\mathbf{M}) = 4619$), ensuring robustness under numerical factorization.

\subsection{Thin SVD-Truncated Matrix Projection Architecture}
\label{sec:methods:svd}

Reconstructing high-fidelity profiles from a sparse tower array ($M = 8$ levels) onto an $N=32$ spectral expansion is an inherently rank-deficient inverse problem. Evaluating the basis at instrument elevations yields the rectangular design matrix $\mathbf{H} \in \mathbb{R}^{M \times (N+1)}$:
\begin{equation}
H_{ij} = T_{j-1}(\xi(z_i))
\label{eq:h_matrix_def}
\end{equation}

To avoid unstable normal-equation inversion and unnecessary parameterized regularization, we use the economy (thin) singular value decomposition (SVD):
\begin{equation}
\mathbf{H} = \mathbf{U}_r \mathbf{S}_r \mathbf{V}_r^T
\label{eq:thin_svd}
\end{equation}
where $\mathbf{U}_r \in \mathbb{R}^{8 \times 8}$ and $\mathbf{V}_r \in \mathbb{R}^{33 \times 8}$ have orthonormal columns, and $\mathbf{S}_r = \operatorname{diag}(s_1, s_2, \ldots, s_8)$ contains the ordered singular values. The thin form restricts all operations to the data-constrained subspace and avoids explicit manipulation of null-space modes.

To separate resolvable physical modes from machine-level noise, we define a dynamic truncation tolerance $\tau$:
\begin{equation}
\tau = \max\left(M, N+1\right) \cdot s_1 \cdot \epsilon_{\mathrm{mach}}
\label{eq:dynamic_tolerance}
\end{equation}
where $\epsilon_{\mathrm{mach}} \approx 2.22 \times 10^{-16}$. The effective rank is
\begin{equation}
r_{\mathrm{eff}} = \#\{s_i > \tau\}.
\label{eq:effective_rank}
\end{equation}
Equivalently, the normalized singular-value cutoff is $\lambda_{\mathrm{min}} = \tau / s_1 = \max(M, N+1)\,\epsilon_{\mathrm{mach}}$. The truncated least-squares estimate is then
\begin{equation}
\mathbf{c} = \sum_{i=1}^{r_{\mathrm{eff}}} \frac{\mathbf{u}_i^{\mathsf{T}}\,\mathbf{a}_{\mathrm{obs}}}{s_i}\,\mathbf{v}_i,
\label{eq:svd_reconstruction}
\end{equation}
which is the Moore--Penrose solution restricted to retained singular directions. This guarantees that inferred coefficients remain in the data-supported subspace while unresolved high-order components are suppressed.

\subsection{Smooth Scale Partitioning and Energy Non-Conservation Caveat}
\label{sec:methods:partition}

To separate multi-scale structures without inducing non-local Gibbs ringing or artificial step discontinuities, we deploy a partition-of-unity spectral windowing framework. The global flow field is decomposed into mean synoptic ($\psi_M$), internal gravity wave ($\psi_W$), and micro-turbulent residual ($\psi_T$) scales using coupled hyperbolic tangent transition channels:
\begin{align}
  \psi_M(n) &= \frac{1}{2}\left[1 - \tanh\left(\frac{n - n_M}{\Delta}\right)\right] \label{eq:psi_m_def} \\
  \psi_W(n) &= \frac{1}{2}\left[1 + \tanh\left(\frac{n - n_M}{\Delta}\right)\right] \cdot \frac{1}{2}\left[1 - \tanh\left(\frac{n - n_W}{\Delta}\right)\right] \label{eq:psi_w_def} \\
  \psi_T(n) &= 1 - \psi_M(n) - \psi_W(n) \label{eq:psi_t_def}
\end{align}
where the scale boundaries are set to $n_M = 3$ and $n_W = 12$ with a smooth transition width $\Delta = 1.2$. By design, the partition operators sum identically to unity across the entirety of the spectral spectrum:
\begin{equation}
\psi_M(n) + \psi_W(n) + \psi_T(n) = 1 \quad \forall n \in [0, N]
\label{eq:partition_unity}
\end{equation}

Crucially, because these smooth partition vectors are applied as coefficient multipliers rather than projections onto true orthogonal eigenvectors of the mass matrix $\mathbf{M}$, \emph{global physical energy content is not strictly additive across the partitioned sub-domains}. When evaluating the individual scale components ($\mathbf{c}^{(M)} = \operatorname{diag}(\psi_M)\mathbf{c}$, $\mathbf{c}^{(W)} = \operatorname{diag}(\psi_W)\mathbf{c}$, and $\mathbf{c}^{(T)} = \operatorname{diag}(\psi_T)\mathbf{c}$), the cross-product metrics do not vanish identically:
\begin{equation}
\langle \mathbf{c}, \mathbf{M} \mathbf{c} \rangle = \sum_{j \in \{M,W,T\}} \langle \mathbf{c}^{(j)}, \mathbf{M} \mathbf{c}^{(j)} \rangle + \sum_{j \neq k} \langle \mathbf{c}^{(j)}, \mathbf{M} \mathbf{c}^{(k)} \rangle
\label{eq:energy_overlap}
\end{equation}
The off-diagonal cross-terms $\langle \mathbf{c}^{(j)}, \mathbf{M} \mathbf{c}^{(k)} \rangle \neq 0$ physically represent spatial and spectral scale overlap regions within the boundary layer transitions. While this cross-term behavior accurately mirrors the coupled wave-turbulence interaction concepts emphasized by \citet{Sun2015}, total energy balance estimates derived via this windowing scheme must explicitly track these non-orthogonal interaction components rather than assuming direct linear superposition.

\subsection{Information-Theoretic Diagnostic Features}
\label{sec:methods:diagnostics}

\subsubsection{Effective Modal Dimension ($D_{\mathrm{eff}}$) with Minimum Energy Filtering}

We compute the effective modal dimension $D_{\mathrm{eff}}$ as the exponential of the spectral Shannon entropy, converting a raw information metric into an interpretable value representing the number of actively participating, equiprobable spectral modes:
\begin{equation}
D_{\mathrm{eff}} = \exp\left( -\sum_n p_n \log p_n \right),
\label{eq:d_eff_def}
\end{equation}
where $p_n = c_n^2 / \sum_m c_m^2$ is the normalized energy in mode $n$.

Because $D_{\mathrm{eff}}$ operates as a scale-invariant shape diagnostic, a quiescent profile under weak background conditions could mathematically register an artificially inflated modal dimension due to numerical precision variations. To prevent low-amplitude instrumental noise from triggering false high-dimensional turbulence classifications, we introduce a regularized minimum energy masking operator before computing the spectral probabilities:
\begin{equation}
c_n^2 \leftarrow \begin{cases}
c_n^2, & \text{if } \sum_{m=0}^{N} c_m^2 \ge \mathcal{E}_{\text{floor}} \\
0, & \text{otherwise}
\end{cases}
\label{eq:energy_floor}
\end{equation}
where the threshold floor is defined dynamically based on historical instrument variance limits, fixed at $\mathcal{E}_{\text{floor}} = 10^{-4} \cdot \max_{t} \left(\|\mathbf{a}_{\mathrm{obs}}\|^2\right)$. This restriction forces highly quiescent, laminar configurations cleanly down to $D_{\mathrm{eff}} \to 1.0$.

\subsubsection{Wave Energy Fraction ($F_W$) and Spectral Curvature ($\chi_N$)}

The wave energy fraction isolates metric-weighted energy in the wave band relative to total energy:
\begin{equation}
F_W = \frac{\langle \mathbf{c}^{(W)}, \mathbf{M} \mathbf{c}^{(W)} \rangle}{\langle \mathbf{c}, \mathbf{M} \mathbf{c} \rangle}
\label{eq:f_w}
\end{equation}

Spectral curvature ($\chi_N$) measures the concentration of energy in high-frequency modes relative to low-frequency modes, providing a proxy for structural sharpness:
\begin{equation}
\chi_N = \frac{\sum_{n=2}^{N} (n/N)^2 \, c_n^2}{\sum_{n=0}^{N} (n/N)^2 \, c_n^2}
\label{eq:chi_n}
\end{equation}

% ============================================================================
\section{Theoretical Extensions and Model Interfacing}
\label{sec:extensions}

\subsection{Physical Linking to the Ozmidov Buoyancy Scale}
\label{sec:extensions:ozmidov}

To transition the statistical boundaries identified by the Gaussian Mixture Model from pure observational diagnostics into formal physical constraints, we map the relationship between the effective modal dimension ($D_{\mathrm{eff}}$) and the Ozmidov length scale ($L_O$). The Ozmidov scale defines the maximum vertical isotropic eddy diameter possible within a stably stratified fluid before buoyancy forces suppress overturning motions:
\begin{equation}
L_O = \sqrt{\frac{\epsilon}{N_b^3}}
\label{eq:ozmidov_def}
\end{equation}
where $\epsilon$ represents the local turbulent TKE dissipation rate, and $N_b = \sqrt{\frac{g}{\theta_{\mathrm{ref}}} \frac{\partial \theta}{\partial z}}$ is the Brunt-Väisälä buoyancy frequency.

Within our Riemannian manifold space, the active spectral degrees of freedom correspond directly to the physical capability of the boundary layer to sustain multi-scale cascade structures. Let $\Delta z_{\text{local}} = J(\xi) \Delta \xi$ represent the local spatial grid resolution at height $z$. We state that the effective modal dimension undergoes an informational collapse when the buoyancy forces constrain the physical eddy size below the structural features resolvable by the dominant basis functions:
\begin{equation}
D_{\mathrm{eff}} \propto \mathcal{F}\left(\frac{L_O}{\Delta z_{\text{local}}}\right) = \gamma_{\text{scal}} \left( \frac{\epsilon}{N_b^3 \cdot [J(\xi)\Delta \xi]^2} \right)^{1/3}
\label{eq:d_eff_ozmidov_link}
\end{equation}
where $\gamma_{\text{scal}}$ is a non-dimensional structural scaling constant.

Under Continuous Turbulence conditions (Cluster 1), $L_O \gg \Delta z_{\text{local}}$, allowing unconstrained three-dimensional energy transfers that populate the entire spectral matrix ($D_{\mathrm{eff}} > 20$). Conversely, within the Wave-Dominated state (Cluster 2), intense radiative cooling drives $N_b \to \text{max}$ and collapses $\epsilon \to 0$, forcing $L_O \ll \Delta z_{\text{local}}$. This spatial restriction suppresses higher-order turbulent basis modes, causing the physical system to contract onto a low-rank, highly compressed gravity-wave manifold baseline ($D_{\mathrm{eff}} \sim 4$--$6$).

\subsection{Resolution of Grid-Topology Mismatch via Virtual Towers}
\label{sec:extensions:virtual_tower}

A critical barrier to utilizing tower-based spectral diagnostics for the validation of mesoscale weather prediction models (e.g., WRF) is the structural grid-topology mismatch. Numerical simulations utilize terrain-following hydrostatic pressure coordinates ($\eta$), distributing prognostic variables across highly coarse, vertically expanding grid blocks that cannot directly interface with a localized, ultra-dense hyperbolic Chebyshev workspace.

To bridge this structural divide, we define a generalized ``Virtual Tower'' interpolation operator $\mathcal{P}_{\text{VT}}: \Omega_{\eta} \to \Omega_{\xi}$. Let $\mathbf{\Phi}_{\eta} = [\phi(1), \phi(2), \dots, \phi(K_{\max})]^T$ represent a discrete variable vector evaluated on the native model grid column at height coordinates $\mathbf{z}_{\eta}$. We map this unaligned state directly into the continuous spectral coefficient space $\mathbf{c}$ via a metric-consistent minimization process:
\begin{equation}
\mathbf{c} = \arg\min_{\mathbf{c}^*} \sum_{k=1}^{K_{\max}} \mathcal{W}*k \left| \phi(\eta_k) - \sum*{n=0}^{N} c_n^* T_n\left(\mathcal{T}^{-1}(z(\eta_k))\right) \right|^2
\label{eq:virtual_tower_opt}
\end{equation}
where $\mathcal{W}_k = (\Delta z_{\eta, k})^{-1}$ represents a local spatial grid-weighting operator, and $\mathcal{T}^{-1}(z)$ represents the analytical physical-to-computational mapping derived by inverting Eq.~\ref{eq:hyperbolic_stretch}:
\begin{equation}
\xi(z) = \frac{(1 - \alpha_s^2)(z - z_{\min}) [1 + 2\alpha_s] - (z_{\max} - z_{\min})(1 + \alpha_s)}{(z_{\max} - z_{\min})(1 + \alpha_s) - \alpha_s(1 - \alpha_s^2)(z - z_{\min})}
\label{eq:inverse_stretch}
\end{equation}

By solving Eq.~\ref{eq:virtual_tower_opt} at every model time-step using a pre-computed pseudo-inverse of the unaligned grid evaluation matrix, numerical models can output real-time synthetic trajectories of $D_{\mathrm{eff}}(t)$ and $F_W(t)$. This allows direct, apples-to-apples validation against physical tower instruments without destroying the structural geometric features of the inversion layers.

\subsection{Spectral Anti-Aliasing for Forward Predictive Closures}
\label{sec:extensions:antialias}

If the metric-consistent decomposition framework is utilized as a active basis for forward-predictive sub-mesoscale forecasting equations, the non-linear convective acceleration terms ($\mathbf{u} \cdot \nabla \mathbf{u}$) will inevitably trigger spectral aliasing. Non-linear cross-multiplication of Chebyshev coefficients drives energy into modes exceeding the maximum resolved index $N$, which then fold back into low-frequency channels as artificial numerical noise, destabilizing the time-integration scheme.

To enforce mathematical stability when projecting non-linear structural couplings, we implement a modified pseudospectral $2/3$-rule anti-aliasing truncation operator $\mathcal{A}$. Let $\mathbf{c}^{(f \cdot g)}$ represent the unfiltered spectral coefficient array resulting from the collocation multiplication of two arbitrary profile fields $f(z)$ and $g(z)$. The anti-aliasing operator truncates high-frequency spectral contributions before they can pollute the metric workspace:
\begin{equation}
\mathcal{A}\left(c_n^{(f \cdot g)}\right) = \begin{cases}
c_n^{(f \cdot g)}, & \text{if } n \le \lfloor \frac{2}{3}N \rfloor \\
0, & \text{if } n > \lfloor \frac{2}{3}N \rfloor
\end{cases}
\label{eq:two_thirds_rule}
\end{equation}

At our operational resolution of $N=32$, the anti-aliasing filter securely cuts off modes above $n_{\text{alias}} = 21$. This ensures that all resolved multi-scale forward-predictive interactions reside safely within the well-conditioned data-constrained subspace, permitting long-term time integration of stable boundary layer microfront transitions without numerical blow-up.

% ============================================================================
\section{Results}
\label{sec:results}

\subsection{Manifold Architecture Validation}
\label{sec:results:manifold}

The UnifiedManifoldWorkspace with $N = 32$ modes on the CASES-99 domain achieved minimum node spacing $\Delta z_{\min} = \SI{2.85}{\milli\meter}$ at $z = \SI{1.5}{\meter}$, maximum node spacing $\Delta z_{\max} = \SI{9.43}{\meter}$ at $z \approx \SI{50}{\meter}$, and a mass matrix condition number $\kappa(\mathbf{M}) = 4619$.

Figure \ref{fig:grid} details the structural composition of the computational grid configuration and basis function derivatives. The rapid, metric-consistent compression focused near the surface captures the acute thermal gradient structure of SBL inversions without introducing numerical instabilities at the boundaries.

\begin{figure}[htbp]
\centering
\includegraphics[width=0.95\linewidth]{manifold_geometry_plots.pdf}
\caption{(a) Hyperbolic grid mapping showing node density concentration at the surface ($z = 1.5$ m) with compactification parameter $\alpha_s = 0.05$. (b) Derivatives of Chebyshev basis polynomials $dT_n/dz$ showing spectral sensitivity at multiple scales.}
\label{fig:grid}
\end{figure}

\subsection{SVD Rank Efficiency Across Physical Regimes}
\label{sec:results:svd}

Across 1386 valid CASES-99 profiles, the system maintained a mean effective rank of $\bar{r}_{\mathrm{eff}} = 7.2 \pm 1.1$. Crucially, evaluating the rank metrics across the distinct GMM-identified clusters revealed a striking physical pattern:
\begin{itemize}
  \item \textbf{Continuous Turbulence (Cluster 1)}: Mean $r_{\mathrm{eff}} = 7.9 \pm 0.2$. All 8 physical observations remain fully resolved and unconstrained, filling the structural matrix space.
  \item \textbf{Wave-Dominated (Cluster 2)}: Mean $r_{\mathrm{eff}} = 5.1 \pm 1.4$. The singular spectrum drops sharply, compressing onto 4 to 6 dominant singular values.
  \item \textbf{Intermittent Shear (Cluster 3)}: Mean $r_{\mathrm{eff}} = 7.3 \pm 0.8$. Shows standard transitional behavior across the mathematical boundaries.
\end{itemize}

This rank compression ($\Delta_{\mathrm{rank}} = 24$ relative to the full 33 spectral dimensions) provides quantitative proof that wave-dominated stable layers natively occupy a compressed lower-dimensional manifold. This confirms that low-rank behavior is an intrinsic structural property of the physical fluid system and not an artifact of arbitrary basis function truncation.

\subsection{Visual Regime Validation planes}
\label{sec:results:tier1}

\subsubsection{Plane 1: Energy--Dimension Phase Space ($D_{\mathrm{eff}}$ vs. $F_W$)}

Figure \ref{fig:tier1_plane1} displays the scatter of all 1386 profiles in the $(D_{\mathrm{eff}}, F_W)$ phase space plane. Three highly separated clusters emerge, achieving an aggregate Silhouette verification score of $\bar{S} = 0.582$ for $K=3$. Mean statistical properties for each regime are compiled in Table \ref{tab:regime_stats}.

\begin{table}[h]
\centering
\caption{Regime Cluster Statistics from the Metric-Consistent GMM Processing}
\label{tab:regime_stats}
\begin{tabular}{l|rrrr|rr}
\hline
Regime & $\langle D_{\mathrm{eff}} \rangle$ & $\langle F_W \rangle$ & $\langle \mathrm{Ri}_g \rangle$ & $\langle \chi_N \rangle$ & Count & Frac. \\
\hline

1. Continuous Turb.   & $23.4 \pm 4.1$ & $0.08 \pm 0.06$ & $0.12 \pm 0.15$ & $0.18 \pm 0.12$ & 412 & 0.297 \\
2. Wave-Dominated     & $5.2 \pm 1.8$  & $0.81 \pm 0.11$ & $0.38 \pm 0.18$ & $0.26 \pm 0.14$ & 389 & 0.281 \\
3. Intermittent Shear & $10.6 \pm 4.2$ & $0.31 \pm 0.15$ & $0.22 \pm 0.16$ & $0.32 \pm 0.16$ & 585 & 0.422 \\
\hline
\end{tabular}
\end{table}

\begin{figure}[htbp]
\centering
\includegraphics[width=0.95\linewidth]{tier1_plane1.pdf}
\caption{Scatter of all 1386 CASES-99 profiles in the effective modal dimension--wave energy fraction plane $(D_{\mathrm{eff}}, F_W)$. Cluster boundaries: Continuous Turbulence (red), Wave-Dominated (blue), and Intermittent Shear Bursts (yellow).}
\label{fig:tier1_plane1}
\end{figure}

The physical interpretation of these distinct boundaries maps directly to the flow topology:
\begin{itemize}
  \item \textbf{Cluster 1 (Continuous Turbulence)}: Localized exclusively at high dimensions ($D_{\mathrm{eff}} > 20$) and suppressed wave signals ($F_W < 0.15$). This matches unconstrained, shear-driven boundary layer mixing.
  \item \textbf{Cluster 2 (Wave-Dominated)}: Tightly bounded at $D_{\mathrm{eff}} \sim 5$ and elevated wave dominance ($F_W > 0.75$). Proved via Section \S\ref{sec:extensions:ozmidov} to represent buoyancy-collapsed configurations tracking macro-scale internal gravity waves.
  \item \textbf{Cluster 3 (Intermittent Shear Bursts)}: Occupies the middle ground ($D_{\mathrm{eff}} \in [8, 15]$, $F_W \in [0.2, 0.5]$), acting as the active fluid transition state.
\end{itemize}

\subsubsection{Plane 2: Gradient Stability vs. Structural Profile Geometry}

To map how traditional physical parameter domains interface with our derived manifold features, we construct a second diagnostic subspace using the gradient Richardson number ($\mathrm{Ri}_g$) and high-frequency spectral curvature ($\chi_N$). This validation plane (Figure \ref{fig:tier1_plane2}) isolates local thermodynamic stability boundaries directly from non-local geometric features.

\begin{figure}[htbp]
\centering
\includegraphics[width=0.95\linewidth]{tier1_plane2.pdf}
\caption{Scatter of the 1386 CASES-99 profiles in the gradient stability--spectral curvature plane $(\mathrm{Ri}_g, \chi_N)$. Cluster color coding is preserved from Plane 1. The classical critical turbulence threshold ($\mathrm{Ri}_c = 0.25$) is marked as a dashed horizontal reference line.}
\label{fig:tier1_plane2}
\end{figure}

The spatial distribution within the $(\mathrm{Ri}_g, \chi_N)$ plane reveals key structural features: Continuous Turbulence maps safely within the weakly stratified shear zone below the critical limit ($\langle \mathrm{Ri}_g \rangle = 0.12$), while the Wave-Dominated state extends far beyond the classical laminar threshold ($\langle \mathrm{Ri}_g \rangle = 0.38$). Crucially, Intermittent Shear Bursts exhibit a uniquely elevated spectral curvature signature ($\langle \chi_N \rangle = 0.32 \pm 0.16$), tracking the formation of localized microfront steps and density discontinuities. Because these microfront features are extremely narrow, their intense local shear is washed out by spatial averaging across tower tiers, causing standard gradient metrics to fail while our spectral curvature successfully captures the structural transition.

\subsection{Multi-Day Synoptic Temporal Evolution and Invariance}
\label{sec:results:synoptic}

To confirm that these three regimes track coherent physical events across the timeline of the campaign rather than fluctuating as random numerical noise, we reconstruct the continuous temporal evolution of the boundary layer across the 10-day intensive observation period (IOP) window.

\begin{figure}[htbp]
\centering
\includegraphics[width=0.95\linewidth]{campaign_synoptic_evolution.pdf}
\caption{Continuous 240-hour timeline tracking 10-minute-averaged spectral features across the CASES-99 intensive observation window (22--31 October 1999). Top panel: Effective modal dimension $D_{\mathrm{eff}}$ trajectory. Bottom panel: Wave energy fraction $F_W$ tracking curve.}
\label{fig:synoptic_timeline}
\end{figure}

The multi-day timeline demonstrates remarkable structural regularity. Each evening transition (typically around 17:30 to 19:00 local time) initiates with a rapid collapse of daytime convective mixing, dropping $D_{\mathrm{eff}}$ from unconstrained high dimensions directly down into the smooth, shear-driven profile matrix of Cluster 1. As radiative surface cooling intensifies toward midnight, generating deep nocturnal temperature inversions, the background flow frequently decouples from the surface friction layer. This decoupling triggers a clean, structural transition into the highly compressed Wave-Dominated state (Cluster 2), where the wave energy fraction surges past $F_W > 0.80$ and $D_{\mathrm{eff}}$ stabilizes near its low-rank baseline of $\sim 5.2$. The long-term preservation of these cluster boundaries proves that our manifold coordinates map invariant physical states of the nocturnal boundary layer.

\subsection{Sponge-Layer Boundary Verification Test}
\label{sec:results:sponge}

A primary vulnerability of high-order spectral projections in sparse observational columns is the generation of spurious boundary reflections. If extracted sub-mesoscale gravity waves encounter an unphysically rigid upper boundary condition, artificial energy reflection can propagate downward, corrupting lower-level flux estimates and inflating the apparent intermittency signature. To verify that our smooth scale-partition functions ($\psi_M, \psi_W, \psi_T$) act as an effective numerical sponge layer, we execute a synthetic wave propagation benchmark test.

We introduce a synthetic, monochromatic internal gravity wave with an upward group velocity into the computational workspace at the lower tower levels:
\begin{equation}
w_{\mathrm{syn}}(z, t) = A_0 \exp(i [k_x x + m z - \omega t])
\label{eq:wave_syn}
\end{equation}
where the vertical wavenumber is set to match a characteristic nocturnal wave structure ($m = 0.25\,\text{m}^{-1}$), and the initial amplitude is fixed at $A_0 = 0.5\,\text{m s}^{-1}$. We then measure the wave reflection coefficient $\mathcal{R} = |A_{\mathrm{reflected}} / A_{\mathrm{incident}}|$ at the upper computational boundary ($z_{\max} = 55.0\,\text{m}$) across different grid stretching and partitioning frameworks.

\begin{figure}[htbp]
\centering
\includegraphics[width=0.85\linewidth]{wave_reflection_test.pdf}
\caption{Sponge-layer wave reflection verification test tracking spatial decay of the normalized upward-propagating synthetic gravity wave amplitude as it approaches the upper boundary of the computational workspace ($z \to 55$ m).}
\label{fig:sponge_test}
\end{figure}

The benchmark test confirms that the unified manifold framework successfully dampens non-physical wave reflections. When an unweighted, standard Chebyshev projection is applied to this sparse setup, a high reflection coefficient ($\mathcal{R}_{\text{standard}} = 0.124$) is observed, creating clear standing-wave artifacts across the lower observation heights. In contrast, our optimized baseline configuration—combining hyperbolic compactification ($\alpha_s = 0.05$) with the smooth partition-of-unity spectral masks—achieves an exceptional reflection coefficient of $\mathcal{R}_{\text{optimized}} = 3.72 \times 10^{-6}$, representing an approximate $3.3 \times 10^4$ reduction in boundary-induced noise amplification. This exceptional attenuation confirms that the spectral windowing functions safely absorb high-frequency grid-scale modes before they can project onto the upper boundary nodes, ensuring that the low-rank states observed in the CASES-99 campaign data reflect genuine physical structures.

% ============================================================================
\section{Discussion and Conclusions}
\label{sec:conclusions}

By projecting sparse meteorological tower observations onto a metric-consistent Riemannian manifold, this study has resolved a long-standing diagnostic bottleneck in stable boundary layer meteorology. The application of our pseudospectral architecture to the CASES-99 intensive observation period has demonstrated that the traditional, highly variable continuum of nocturnal stable flows can be cleanly separated into three statistically distinct, invariant physical regimes without relying on empirical flux-profile look-up tables.

The numerical success of this framework centers on the correction of the Chebyshev coordinate mapping baseline within the \emph{UnifiedManifoldWorkspace}. Resolving the array-ordering inversion in the node distribution ensured that the observed instrument heights ($1.5\,\text{m}$ to $55.0\,\text{m}$) were accurately mapped across the computational domain $\xi \in [-1, 1]$, preventing the structural rank collapse that previously forced metrics into an uninformative rank-1 subspace. With the physical variance profile restored, the SVD singular spectrum cleanly identified that Wave-Dominated states (Cluster 2) natively contract into a compressed low-rank subspace ($D_{\mathrm{eff}} = 5.2 \pm 1.8$). This structural rank preservation provides solid evidence that the observed low-dimensional configurations represent true physical states of the atmospheric fluid system, rather than artifacts of mathematical truncation.

Furthermore, the linear independence demonstrated within the feature orthogonality matrix offers a path forward for overriding the current parameterization failures of Monin--Obukhov similarity theory (MOST). Because our high-frequency spectral curvature parameter ($\chi_N$) exhibits near-zero correlation with localized stability indicators like the gradient Richardson number ($\mathrm{Ri}_g$), it provides a separate, geometric diagnostic axis capable of tracking narrow microfront steps and internal gravity wave steepening. Incorporating $D_{\mathrm{eff}}$ and $\chi_N$ as multi-scale state indicators within next-generation mesoscale models would allow parameterizations to dynamically adjust local mixing coefficients based on global manifold geometries, capturing the physics of sudden intermittent turbulent bursts that standard single-column models miss.

Ultimately, the combination of non-uniform hyperbolic compactification and smooth partition-of-unity spectral windowing provides a mathematically rigorous, boundary-consistent diagnostic suite. The complete elimination of unphysical upper-boundary wave reflections ($\mathcal{R} \sim 3.7 \times 10^{-6}$) ensures that this pipeline can be safely exported to other dense micrometeorological tower domains (e.g., SHEBA, FLUXNET). This enables physically interpretable, automated regime classification with direct utility for data assimilation, turbulence closure development, and high-fidelity mesoscale model validation.

% ============================================================================
\section*{Acknowledgments}
The author expresses sincere gratitude to the National Center for Atmospheric Research (NCAR) Earth Observing Laboratory (EOL) for maintaining and providing the public data repository for the CASES-99 campaign. We also acknowledge the invaluable open-source contributions of the Julia language ecosystem, particularly the \texttt{LinearAlgebra.jl}, \texttt{DataFrames.jl}, and \texttt{CSV.jl} implementation libraries that powered this parallel pipeline engine.

\bibliographystyle{apsrev4-2}
\bibliography{references}

\end{document}
